\section{Methodology}
\label{sec:methodology}

Our methodology follows five steps aligned with the research objectives:
(1)~standardise the datasets and reproduce response-only KT baselines
(RO1); (2)~implement the proposed knowledge-component (KC) /
code-structure-aware variant (RO2); (3)~measure prediction accuracy under
student-level cross-validation (RO3); (4)~measure downstream sequencing
utility and its correlation with prediction (RO4); and (5)~test
cross-dataset and cross-domain transfer (RO5). The selected model and
policy are the ones the AlgoVenture adaptive engine adopts for mastery
estimation (UC6) and next-exercise sequencing (UC7).

% =====================================================================
\subsection{Problem Formulation}
\label{sec:formulation}
We keep only the one equation needed to state precisely what a KT model
predicts and where the programming signals enter. For a learner, the
interaction sequence is $\{(q_t, c_t, r_t, s_t)\}_{t=1}^{T}$: at step $t$
the learner attempts problem $q_t$, tagged with knowledge components
$c_t$, with response $r_t \in \{0,1\}$ (1 = correct) and submitted source
code $s_t$. The task is to predict the next response,
\begin{equation}
\hat{r}_{t+1} = P\big(r_{t+1}=1 \mid q_{t+1}, c_{t+1},\,
                       \text{history}_{1:t}\big).
\label{eq:kt-task}
\end{equation}
\emph{What this solves:} it fixes the prediction target used by every
model in the study, so baselines and our variant are compared on exactly
the same quantity. Response-only baselines use only $(q,c,r)$ in the
history; our variant additionally uses the code $s$.

% =====================================================================
\subsection{Pipeline and Baselines (RO1)}
\label{sec:pipeline}
We build a reproducible pipeline on top of the pyKT
toolkit~\cite{liu2022pykt}, which provides standardised preprocessing and
reference implementations, so that measured differences reflect the model
rather than the data plumbing. The response-only baselines are
DKT~\cite{piech2015dkt}, SAKT~\cite{pandey2019sakt},
AKT~\cite{ghosh2020akt}, and simpleKT~\cite{liu2023simplekt};
BKT~\cite{corbett1994bkt} is a classical reference. This baseline set is
also the promotion gate for the product: a KT model is only adopted by
the app after being evaluated against these baselines on AUC/RMSE.

% =====================================================================
\subsection{Proposed KC / Code-Structure-Aware Model (RO2)}
\label{sec:model}
We extend a strong attention baseline (AKT or simpleKT) with two extra
input channels, described in words rather than layer-by-layer equations:

\begin{itemize}
  \item \textbf{Knowledge-component channel:} an embedding of the KC
        tag(s) of the current problem. \emph{Why:} it lets the model share
        evidence across problems that exercise the same concept, instead
        of treating each problem in isolation.
  \item \textbf{Code-structure channel:} a fixed-length feature vector
        built from the abstract syntax tree (AST) of the submitted code
        (counts and nesting depths of loops, conditionals, calls, etc.),
        or a learned code embedding. \emph{Why:} two submissions can both
        be ``wrong'' yet reveal very different misconceptions; code
        structure exposes that signal, which a correct/incorrect label
        hides.
\end{itemize}

The two channels are combined with the base model's interaction
embedding by a simple fusion (concatenation with a linear projection, or
a gated sum), and the base attention model then predicts
$\hat{r}_{t+1}$ as in Eq.~\eqref{eq:kt-task}. We deliberately keep the
fusion simple so that any measured gain is attributable to the
\emph{signals} (KC, code) rather than to a more complex predictor.

\paragraph{Ablations.}
To isolate each signal we run four configurations: (i)~response-only;
(ii)~response~+~KC; (iii)~response~+~code; and (iv)~the full model
(response~+~KC~+~code). Comparing (ii) and (iii) against (i) isolates each
signal; comparing (iv) against them tests whether the signals are
complementary. This is the direct evidence for hypothesis~H1a.

% =====================================================================
\subsection{Prediction Evaluation (RO3)}
\label{sec:pred-eval}
We report AUC, RMSE, and accuracy under 5-fold \emph{student-level}
cross-validation, so no student appears in both training and test folds
(guarding against the label-leakage pitfall discussed in
Section~\ref{sec:threats}). We report mean and standard deviation across
folds and test significance with paired tests plus effect sizes
(hypotheses~H1, H1a).

% =====================================================================
\subsection{Sequencing-Utility Evaluation (RO4)}
\label{sec:sequencing}
This is the core of the study: does a more accurate predictor actually
sequence exercises better? We evaluate sequencing offline, using a
trained KT model as a learner simulator. A policy $\pi$ repeatedly picks
the next problem, the simulator returns a predicted response, and the
state updates. We summarise a run by the utility $U$---the expected
mastery gain over a horizon of $H$ recommendations:
\begin{equation}
U(\pi) = \mathbb{E}\big[\, m_{H} - m_{0} \,\big],
\label{eq:utility}
\end{equation}
where $m_0$ and $m_H$ are estimated mastery before and after the
sequence. \emph{What this solves:} it turns ``good teaching'' into one
comparable number per model, so we can line utility up against accuracy.
Algorithm~\ref{alg:seq} states the procedure.

\begin{algorithm}[t]
\caption{Offline sequencing-utility evaluation}
\label{alg:seq}
\begin{algorithmic}[1]
\STATE \textbf{Input:} KT simulator $M$, policy $\pi$, horizon $H$,
       start states $\mathcal{S}$
\STATE $U \gets 0$
\FORALL{start state $z_0 \in \mathcal{S}$}
  \STATE $m_0 \gets \textsc{Mastery}(M, z_0)$;\quad $z \gets z_0$
  \FOR{$h = 1$ \TO $H$}
    \STATE $q \gets \pi(z)$;\quad $r \gets M(z, q)$;\quad
           $z \gets \textsc{Update}(z, q, r)$
  \ENDFOR
  \STATE $U \gets U + \big(\textsc{Mastery}(M, z) - m_0\big)$
\ENDFOR
\STATE \textbf{return} $U / |\mathcal{S}|$
\end{algorithmic}
\end{algorithm}

We then correlate each model's predictive accuracy with its sequencing
utility (Pearson and Spearman, with a bootstrap confidence interval),
which directly answers research question~(b) and tests hypothesis~H3.
Because a KT simulator is an imperfect proxy for real learners, we add
robustness checks (multiple simulators, sensitivity to $H$) and read
utility as offline evidence, not a deployment guarantee
(Section~\ref{sec:threats}).

% =====================================================================
\subsection{Generalisation Evaluation (RO5)}
\label{sec:gen-eval}
We test cross-dataset transfer (train on one programming dataset, test on
another) and cross-domain transfer (programming vs.\ mathematics),
reporting transfer AUC against the within-dataset reference so the size of
any drop is visible (hypotheses~H4, H5).
